% TOP_PROBLEM: {"id": 30006969, "title": "Spectral Determination of Convex Planar Domains", "category_id": 8, "category": {"id": 8, "name": "analysis", "display_name": "Analysis", "description": "Limits, continuity, calculus, and function theory.", "slug": "analysis", "order_index": 8, "created_at": "2026-07-31T15:26:25.671Z"}, "status": "open"}
% FIRST_POSTED: 2026-09-27
% !TeX program = lualatex
% ATTEMPT_STATUS: unresolved
\documentclass[11pt]{article}
\usepackage[margin=25mm]{geometry}
\usepackage{fontspec}
\setmainfont{Latin Modern Roman}
\usepackage{amsmath,amssymb,mathtools}
\usepackage{unicode-math}
\setmathfont{Latin Modern Math}
\usepackage{xurl,hyperref}
\hypersetup{colorlinks=true,urlcolor=blue}
\setcounter{secnumdepth}{0}
\setlength{\parindent}{0pt}
\setlength{\parskip}{5pt}
\setlength{\emergencystretch}{3em}
\begin{document}
\section*{\texorpdfstring{Spectral Determination of Convex Planar Domains}{Spectral Determination of Convex Planar Domains}}\label{30006969}
\subsection{Definitions and mathematical statement}
For a bounded open convex domain $\Omega\subset{\mathbb{R}}^2$, the Dirichlet Laplacian is the self-adjoint operator associated with the quadratic form $u\mapsto\int_\Omega|\nabla u|^2$ on $H_0^1(\Omega)$. Write its discrete spectrum, with multiplicities, as
\[
 0<\lambda_1(\Omega)\le\lambda_2(\Omega)\le\cdots.
\]
The question is whether
\[
 \lambda_j(\Omega_1)=\lambda_j(\Omega_2)\ (j\ge1)
 \quad\Longrightarrow\quad \Omega_2=U\Omega_1+b
\]
for an orthogonal matrix $U$ and translation $b$. Reflections are included among congruences. The entire spectrum with multiplicities is required; equality of finitely many eigenvalues or heat invariants is weaker. The convexity assumption must be retained when considering examples of isospectral domains.
\par\smallskip\textit{Formulation and concept references:} \href{https://arxiv.org/html/2406.18369v2}{[S1]}.

\subsection{Short English statement}
Does the full Dirichlet spectrum determine a bounded convex planar domain uniquely up to Euclidean congruence?
\subsection{Research attempt}
\paragraph{Scope and source audit.}
This exploratory attempt by GPT-6 Astra Ultra was prepared on 27 September
2026. The target is global uniqueness among all bounded open convex planar
domains for the full Dirichlet spectrum with multiplicities. A statement
restricted to a family, to an infinitesimal deformation, or to finitely many
spectral invariants is a smaller assertion. The supplied survey [S1],
Section 7, identifies the convex-drum question as open. The source check
also found the 2026 manuscript [E2] claiming strictly convex isospectral
examples for the \emph{Steklov} problem. Those concern a different operator
and are not counterexamples to the present statement. No argument from that
manuscript is used below.

The relevant heat-trace input is the two-term Dirichlet asymptotic on
bounded Lipschitz domains, discussed with its precise regularity in [E1].
Bounded convex planar domains with nonempty interior have Lipschitz
boundary, so that input applies even when their boundaries are not smooth.
Higher smooth-boundary coefficients are used only for the explicit smooth
examples below. This distinction prevents a silent smoothness assumption
from narrowing the original question.

\paragraph{Attack map.}
First extract geometric data from the spectrum and solve a restricted
reconstruction problem exactly. Then try to reconstruct an arbitrary convex
body from the area and perimeter. An explicit pair of noncongruent strictly
convex analytic boundaries will disprove that intermediate inference.
Finally retain the complete heat trace instead of its initial coefficients,
and locate the genuinely missing inverse step. These are separate proof and
counterexample tracks: a counterexample to a proposed finite-invariant
criterion is not an isospectral counterexample.

\paragraph{Lemma 492.1: the first two heat coefficients are audible.}
Define
\[
 Z_\Omega(t)=\sum_{j=1}^{\infty}e^{-t\lambda_j(\Omega)},
                 \qquad t>0.
\]
For a bounded Lipschitz planar domain the established heat expansion is
\begin{equation}\label{a492:heat}
 Z_\Omega(t)=\frac{A}{4\pi t}
       -\frac{P}{8\sqrt{\pi t}}+o(t^{-1/2}),\qquad t\downarrow0,
\end{equation}
where $A=|\Omega|$ and $P=\mathcal H^1(\partial\Omega)$. This is the
two-term heat formula, not a pointwise two-term counting law requiring
additional dynamical hypotheses. See [E1], equation (3) and its discussion
of Brown's Lipschitz-boundary theorem. We use that analytic theorem as an
external input, not as a new result proved here.

The extraction itself is immediate but worth making exact:
\begin{equation}\label{a492:extract}
 A=\lim_{t\downarrow0}4\pi t Z_\Omega(t),\qquad
 P=-8\sqrt\pi\lim_{t\downarrow0}\sqrt t
       \left(Z_\Omega(t)-\frac{A}{4\pi t}\right).
\end{equation}
Thus isospectral convex domains have equal area and perimeter. For smooth
simply connected planar domains the next constant coefficient is $1/6$;
the curvature form is $(12\pi)^{-1}\int_{\partial\Omega}\kappa\,ds$,
and the total turning is $2\pi$. The smooth heat formula and the distinction
from corner coefficients are discussed in [S1], Section 4.2.

In the equality case $P^2=4\pi A$, the planar isoperimetric theorem [E3] forces
the domain to be a disk. Hence a disk is determined among the convex
domains in the question by these two audible quantities alone. This uses
the classical equality case of the isoperimetric theorem as an external
geometric input. Away from equality, that argument supplies no uniqueness.

\paragraph{Attempt A: exact reconstruction inside the rectangle family.}
Let $\Omega=(0,a)\times(0,b)$ with $a\ge b>0$. Separation of variables
gives the complete Dirichlet eigenvalue multiset
\begin{equation}\label{a492:rectangle}
 \lambda_{m,n}=\pi^2\left(\frac{m^2}{a^2}+\frac{n^2}{b^2}\right),
                         \qquad m,n\ge1.
\end{equation}
The products of one-dimensional sine bases are a complete orthogonal basis
of $L^2(\Omega)$, so there are no missing eigenfunctions. Put
$\alpha=a^{-2}$ and $\beta=b^{-2}$; then $0<\alpha\le\beta$.
The smallest index pair is $(1,1)$. Among every other pair, either
$m\ge2$, giving a value at least $4\alpha+\beta$, or $n\ge2$,
giving a value at least $\alpha+4\beta\ge4\alpha+\beta$.
Consequently, counting multiplicity,
\[
 \lambda_1=\pi^2(\alpha+\beta),\qquad
 \lambda_2=\pi^2(4\alpha+\beta).
\]
Solving the two linear equations proves
\begin{equation}\label{a492:inverse}
 a=\pi\sqrt{\frac{3}{\lambda_2-\lambda_1}},\qquad
 b=\pi\sqrt{\frac{3}{4\lambda_1-\lambda_2}}.
\end{equation}
For a square the second eigenvalue is double, but its value is still the
one used above. The denominators are positive. The ratio satisfies
$1<\lambda_2/\lambda_1\le5/2$, with equality at $5/2$ exactly for
$a=b$. This proves uniqueness \emph{within the rectangle family} from
just two eigenvalues. It does not show that an arbitrary convex isospectral
competitor must be a rectangle.

There is also a reconstruction from the two heat coefficients: $a+b=P/2$
and $ab=A$, so $a,b$ are the roots of
$z^2-(P/2)z+A=0$. Its discriminant is $(a-b)^2$.
This reveals a stability limitation near a square. Small errors in $A,P$
can produce square-root errors in the recovered side difference. Exact
uniqueness does not by itself provide a uniformly well-conditioned numerical
inverse, and a finite precision eigenvalue match is not an exact spectral
identity.

\paragraph{Attempt B: test whether area and perimeter determine a convex body.}
We now give an explicit obstruction to extending the preceding two-invariant
argument. For an integer $m\ge2$, let
\begin{equation}\label{a492:support}
 h_m(\theta)=R+a_m\cos(m\theta),\qquad
 a_m=\frac{\tau}{\sqrt{m^2-1}},
 \quad 0<\tau<\frac{R}{\sqrt{m^2-1}}.
\end{equation}
With $u=(\cos\theta,\sin\theta)$ and
$u_\perp=(-\sin\theta,\cos\theta)$, define
\[
 r_m(\theta)=h_m(\theta)u(\theta)+h'_m(\theta)u_\perp(\theta).
\]
Its derivative is
\[
 r'_m(\theta)=\rho_m(\theta)u_\perp(\theta),\qquad
 \rho_m=h_m+h''_m=R-\tau\sqrt{m^2-1}\cos(m\theta)>0.
\]
This is a regular analytic strictly convex closed curve with support
function $h_m$. Here is a direct check of the support assertion. For fixed
$\phi$, the derivative of $r_m(\theta)\cdot u(\phi)$ equals
$\rho_m(\theta)\sin(\phi-\theta)$, which is negative on
$(\phi,\phi+\pi)$ and positive on $(\phi+\pi,\phi+2\pi)$.
It therefore has a unique maximum at $\theta=\phi$ modulo $2\pi$,
with value $h_m(\phi)$. Each direction has its unique supporting point;
the curve traces the boundary of its convex hull once and is embedded.
The strictly positive radius of curvature excludes flat segments.

Its perimeter and area follow without a numerical approximation:
\begin{align}
 P_m&=\int_0^{2\pi}\rho_m(\theta)\,d\theta=2\pi R,
                                                       \label{a492:perimeter}\\
 A_m&=\frac12\int_0^{2\pi}\det(r_m,r'_m)\,d\theta
      =\frac12\int_0^{2\pi}(h_m^2-(h'_m)^2)\,d\theta
      =\pi R^2-\frac\pi2\tau^2.                        \label{a492:area}
\end{align}
The integration by parts uses periodicity, and the final equality uses
the zero average of the cosine and the integrals $\int\cos^2(m\theta)
=\int\sin^2(m\theta)=\pi$.

Choose $R=1$, $\tau=1/10$, and $m=2,3$. Both curvature radii are
positive, and the two interiors have identical area $199\pi/200$,
perimeter $2\pi$, and the same constant smooth heat coefficient $1/6$.
They are not congruent. The $m=2$ body is centrally symmetric since
$h_2(\theta+\pi)=h_2(\theta)$. If the $m=3$ body had centre $c$,
its translated support function would be $\pi$-periodic, requiring
\[
       2a_3\cos(3\theta)=2c\cdot u(\theta).
\]
Multiply by $\cos(3\theta)$ and integrate over a period. The right side
is zero by Fourier orthogonality while the left side is $2\pi a_3>0$.
This contradiction excludes every possible centre. Since central symmetry
is preserved by translations, rotations, and reflections, the bodies cannot
be congruent.

This pair refutes reconstruction from the first three displayed heat
coefficients even in the analytic, strictly convex class. It is expressly
\emph{not} asserted to be isospectral. A further geometric diagnostic already
distinguishes them:
\begin{equation}\label{a492:curvature}
 \int_{\partial\Omega_m}\kappa^2\,ds
   =\int_0^{2\pi}\frac{d\theta}{\rho_m(\theta)}
   =\frac{2\pi}{\sqrt{R^2-(m^2-1)\tau^2}}.
\end{equation}
The last integral follows by periodic substitution and the tangent
half-angle substitution in $\int_0^{2\pi}(R-c\cos\theta)^{-1}d\theta$.
This calculation shows that matching the initial data has not even fixed
all elementary curvature statistics. No coefficient of a higher heat term
is needed for the counterexample to the proposed two-invariant method.

\paragraph{Attempt C: retain the full heat trace.}
An equality $Z_{\Omega_1}(t)=Z_{\Omega_2}(t)$ for \emph{every} $t>0$
is equivalent to equality of the complete eigenvalue multisets. To verify
the nontrivial direction, write the lowest distinct eigenvalue as $\mu_1$
with multiplicity $k_1$. Then
\[
 \mu_1=-\lim_{t\to\infty}t^{-1}\log Z(t),\qquad
 k_1=\lim_{t\to\infty}e^{t\mu_1}Z(t).
\]
The remaining tail tends to zero after multiplication by $e^{t\mu_1}$:
for fixed $t_0>0$, dominate its terms for $t\ge t_0$ by the summable
series $e^{-t_0(\lambda_j-\mu_1)}$ and use dominated convergence.
Subtract $k_1e^{-t\mu_1}$ and repeat to recover each successive distinct
eigenvalue and its multiplicity. Discreteness and divergence of the
Dirichlet eigenvalues justify this induction. Thus there is no loss of
information in switching from the whole spectrum to the whole heat trace.

There is a loss in replacing that function by a finite, or merely formal,
small-time expansion. Even two abstract functions whose difference is
$e^{-1/t}$ have identical power asymptotics to every order at zero while
being unequal for $t>0$. This elementary example is not claimed to be a
pair of heat traces of domains. It pinpoints why coefficient matching alone
needs an additional uniqueness theorem before it can substitute for the
full spectral data.

In support coordinates the actual missing implication is now concrete:
for arbitrary convex bodies with support functions $h_1,h_2$, prove that
equality of their Dirichlet heat traces for all positive times forces
\[
 h_2(\theta)=h_1(\varepsilon\theta-\theta_0)+b\cdot u(\theta),
                         \qquad \varepsilon\in\{1,-1\}.
\]
This is precisely the action of orthogonal transformations and translations
on support functions. The formulas \eqref{a492:perimeter}--\eqref{a492:area}
access only two averages of $h$; no argument here extracts its full angular
dependence from $Z(t)$. Positivity of $h+h''$ restricts the class but does
not repair the information loss, as the explicit pair shows.

\paragraph{Verification, boundary cases, and outcome.}
The companion script uses rational arithmetic to enumerate rectangle
eigenvalues $\lambda_{m,n}/\pi^2$ for integer sides from 1 through 12.
For every unordered pair of sides it checks the first two sorted eigenvalues
and the exact inverse formulas, including the square's multiplicity.
It also checks the algebraic curvature-radius positivity for the two support
functions and records numerical quadrature as a diagnostic of the derived
area, perimeter, and curvature integrals. Quadrature is not used to prove
convexity, noncongruence, or equality of any spectrum.

The zero-width rectangle is excluded since the target is an open planar
domain. Reflections are included throughout. The support-function examples
are interiors of compact convex bodies with positive curvature and hence
satisfy all domain hypotheses. The global problem also allows corners,
so an eventual general proof cannot rely solely on their analytic
parametrizations. Equality of a finite number of eigenvalues, or of values
computed to finite precision, would not satisfy its hypothesis.

\textbf{UNRESOLVED.} We have an exact restricted reconstruction theorem,
an attributed disk rigidity argument, a completely explicit failure of
reconstruction from the first heat invariants, and an exact equivalence
between complete heat data and complete spectral data. These do not prove
global uniqueness for convex planar domains. The next substantive tasks are
to extract angular boundary information from spectral quantities beyond the
initial heat coefficients, to prove a rigidity theorem with hypotheses
covering all convex competitors, or to construct two noncongruent convex
domains with an exact full-spectrum identity. The examples here accomplish
none of these final steps and are not proposed counterexamples to the
catalogue conjecture.


\subsection{Sources}
\begin{itemize}
\item[S1]G. Mardby and J. Rowlett, 112 years of listening to Riemannian manifolds, section on planar domains (2024).
\url{https://arxiv.org/html/2406.18369v2}.
\textit{Formulation source.}
\item[E1]Rupert L. Frank and Simon Larson,
\emph{On the error in the two-term Weyl formula for the Dirichlet Laplacian},
arXiv:2001.01876, Introduction, equation (3) and discussion of Brown's
two-term heat formula for Lipschitz domains.
\url{https://arxiv.org/html/2001.01876}.
\textit{External heat-asymptotic input; consulted 27 September 2026.}
\item[E2]Tao Hu, Jiachen Shi, and Quanyu Tang,
\emph{Strictly Convex Steklov-Isospectral Plane Domains},
arXiv:2608.10557v2, 15 August 2026, abstract.
\url{https://arxiv.org/abs/2608.10557v2}.
\textit{Operator-scope check only: this manuscript concerns Steklov,
not Dirichlet, spectra. Its construction is not used here.}
\item[E3]Robert Osserman, \emph{The isoperimetric inequality},
Bulletin of the American Mathematical Society 84 (1978), 1182--1238,
opening statement of the planar inequality and equality case.
\url{https://sites.math.washington.edu/~toro/Courses/20-21/MSF/osserman.pdf}.
\textit{Classical geometric input for the disk subcase; consulted
27 September 2026.}
\end{itemize}
\end{document}
